% !TEX root = ../final-main.tex
\chapter*{Abstract}
\addcontentsline{toc}{chapter}{Abstract}

This report presents the design, fabrication and characterisation of
enhancement-mode GaAs/AlGaAs high-electron-mobility transistors formed by gate
recession. Conventional GaAs/AlGaAs HEMTs are normally-on devices, requiring a
negative gate bias to suppress conduction. The aim of this project was to
investigate whether controlled removal of material beneath the gate could shift
the threshold voltage into the positive region and produce normally-off
operation.

An unrecessed depletion-mode process was first established to provide a
fabrication and electrical baseline. Recessed-gate devices were then developed
through successive fabrication iterations, addressing limitations in recess
geometry, gate layout and etch-depth control. The final non-oxide devices
demonstrated enhancement-mode behaviour, with positive threshold voltages of
approximately \SI{0.8}{\volt} and \SI{1.5}{\volt} for the two recess conditions
investigated. These results show that increasing the recess etch shifts the
threshold voltage further in the positive direction, consistent with reduced
donor charge beneath the gate.

The work demonstrates that gate recession is a viable route to
enhancement-mode operation in the available GaAs/AlGaAs heterostructure. It
also identifies the main practical limits of the process: nanometre-scale etch
control, reduced drain-source current after recession, gate leakage, and the
need to balance oxide passivation against gate control. The final devices
therefore provide both a proof of principle and a basis for assessing how the
process could be improved for more repeatable normally-off GaAs HEMTs.
